\section{Reading a formula without losing the picture}
Sometimes we need to add many quantities. The notation
\[
\sum_{j=1}^{4}a_j=a_1+a_2+a_3+a_4
\]
means add the entries named $a_1$ through $a_4$. The symbol $j$ is a temporary index: it takes the values $1,2,3,4$ in turn. A subscript is a label, not multiplication. For example, $a_2$ is the second entry, whereas $2a$ is twice $a$.

Suppose the entries are $1,3,5,7$. Their total is $16$. Changing the temporary letter from $j$ to $r$ does not change the sum. Changing the upper limit from $4$ to $3$ does: we would obtain $1+3+5=9$. Small pieces of notation can express substantial changes in a claim.

Several collections of numbers have standard names.
\begin{itemize}
\item $\NN=\{0,1,2,\ldots\}$ is the set of \emph{natural numbers}. Braces enclose the members of a collection, and the dots mean that the displayed pattern continues. In this book zero counts as a natural number; when we need a positive number of objects we will say $n\ge1$ explicitly.
\item $\ZZ$ is the set of integers, introduced above.
\item $\QQ$ is the set of \emph{rational numbers}: numbers that can be written as a fraction $a/b$ of integers with $b\ne0$, such as $3/4$, $-5/2$, or $7=7/1$.
\item $\RR$ is the set of \emph{real numbers}: all the numbers on the number line. Every integer and every fraction is a real number, but the number line also contains numbers that are not fractions of integers. Chapter~4 proves that the length $\sqrt2$ of the diagonal of a unit square is one of them.
\end{itemize}
When a letter is said to be real, every point of the number line is an allowed value, including negative numbers, zero, and numbers that are not fractions. The notation $x\ge y$ means $x$ is at least $y$, while $x>y$ means $x$ is strictly larger than $y$. The difference between strict and non-strict inequalities often matters at the boundary of a claim.

The word \emph{positive} means strictly greater than zero. \emph{Nonnegative} means greater than or equal to zero, so it includes zero. Likewise, $x\le y$ means $x$ is at most $y$, and $x<y$ means strictly smaller. A chain such as $0\le x\le1$ states two conditions together: $0\le x$ and $x\le1$.

Later we will use square roots. For a nonnegative real $t$, the symbol $\sqrt t$ denotes the \emph{nonnegative} real number whose square is $t$. For example, $\sqrt9=3$, although both $3$ and $-3$ have square nine. We use the usual real-number fact that this nonnegative square root exists and is unique. The expression $\sqrt{x^2}$ therefore equals the nonnegative magnitude of $x$, not always $x$ itself. At $x=-3$ it equals three. We give this magnitude its own notation in the next chapter.
